\section{Checks of the analytic simplification}
These checks compare the simplified expressions with the original
independent sector results. The algebraic derivations above fix all
coefficients and boundary values before any numerical comparison.

\subsection{Checks of the classical root functions}
The reduction was also checked against the defining integrals at
$X=10^{-6}$, $10^{-2}$, $0.6$, $1$, $10$ and $10^4$, using 70-digit arithmetic.
The substitution $X=4\sinh^2(y/2)$ removes the square-root kernel
from those numerical checks. All 36 endpoint comparisons have absolute
differences below $1.2\times10^{-68}$. Independent differentiation
checks of all six classical expressions at five of these points also
reproduce their defining kernels. These numerical checks supplement
the primitive identities and boundary derivations above. They do not
supply any part of the analytic expressions.

\subsection{Checks of the simplified soft primitive}
At an interior endpoint, the old primitive equals
$\widetilde\Psi+4apf/(\beta E^2dk)$, which permits a direct check
without a numerical integral.  At $z=0.7Z$, evaluating both GPL
expressions with 65 decimal digits gives differences below
$2.6\times10^{-63}$ for $s=1.6$ and $s=10$.  These checks support
the implementation of the exact relation
\eqref{eq:comp-ck-pole-removal}.

\subsection{Checks of the collected soft kernel}
As a cross-check of this reorganization, numerical integration of the
regular kernel gives
\begin{align}
 \mathcal C_K(1.6)&=0.1712265509990548695712562109522182979661\ldots,\\
 \mathcal C_K(10)&=10.0584622326319768377445483326880109081643\ldots.
\end{align}
These values agree with the previous analytic soft expression.
The numerical integrations test the algebraic reorganization and do
not replace the analytic expression.

\subsection{Checks of the soft energy-variable integrals}
Direct quadrature of the defining
$I_1,I_2,I_7$ integrals at $X=10^{-3},0.6,9,100$ agrees with their
energy-variable analytic expressions to better than $3\times10^{-72}$
in absolute value for $I_2$ and $I_7$. The integral evaluations in this
check are not used to determine any coefficient or boundary constant.

\subsection{Checks of the integrated UW expression}
Evaluation of the fixed-alphabet result and the original finite GPL primitive gives
\begin{align} F_{UW}^{\eta}(0.6)&=5.33372142619549103751931327587135840926047\ldots,\\
F_{UW}^{\eta}(9)&=-599.278509012081875840715039026856308771543\ldots.\end{align}
The absolute differences are $1.4\times10^{-65}$ and $3.1\times10^{-63}$, respectively. These checks concern the $UW$ contribution. The conversion to \eqref{eq:comp-uw-integrated-compact} uses the exact coefficient and shuffle identities above.

\subsection{Checks of the matched real contribution}
Fresh evaluations of the original four GPL primitives at $s=1.6$ and $10$ agree with the endpoint formula before matching. Adding the original finite shifts gives agreement with \eqref{eq:comp-real-four-matched-compact} through the 45 significant digits recorded for those original values. These checks use our own integrated expressions. No published total is substituted.

\subsection{Checks of the four-sector endpoints}
The coefficients $A_4,D_4$ extracted from the collected primitive were
also checked against the Laurent expansions of the four original
integrands at $s=1.01$, $1.6$ and $10$, using 70-digit arithmetic.
This checks the endpoint signs and normalization independently of the
primitive coefficient collection.

\subsubsection{The conjugate-pair endpoint}
The function $\mathcal A_{\mathrm{pair}}(s)$ is defined in
\eqref{eq:comp-real-a-pair-endpoint}. Direct quadrature of $P'\ln B/P$ at six points from
$s=1.000001$ to $10^4$ agrees with this expression to an absolute
error below $6\times10^{-70}$ in 70-digit arithmetic.

\subsection{Checks of the direct-box basis conversion}
The independent direct-box conversion is derived in
Section~\ref{sec:comp-virtual-energy-basis}. As a check
of its implementation, evaluation at $s=1.6$ and $s=10$ differs from the
original closed endpoint representation by less than
$3.4\times10^{-78}$.  The exact derivative and boundary arguments,
rather than these two numbers, establish the basis transformation.
